The goal of knowledge graph{\textendash}grounded dialog generation is to generate a dialog response by jointly reasoning over a dialog history and a knowledge graph.
We represent a dialog history as a token sequence, $\TokenSeq = \left[\Token_1, \Token_2, \dots, \Token_{n} \right]$, where $\Token_i \in \Voca$ is the $i$-th token of the dialog history and $\Voca$ denotes the vocabulary set.
A knowledge graph is defined as $\KG = \left(\EntSet, \RelSet, \TripleSet \right)$, where $\EntSet$ is the set of entities and $\RelSet$ is the set of relations.
$\TripleSet$ denotes the set of triplets, $(\Ent_h, \Rel, \Ent_t) \in \TripleSet$, each of which are composed of a head entity $\Ent_h \in \EntSet$, a tail entity $\Ent_t \in \EntSet$, and a relation $\Rel \in \RelSet$ between the two entities.
\jyp{We use $k$-hop subgraph linked to the entities mentioned in the input dialog as retrieval candidates following previous works~\cite{kang2022knowledge}.}
The example of a extracted candidate subgraph is in Figure~\ref{fig:subgraph}.
% \jk{The prev sentence would be confusing to those not famialiry with the approach. Elaborate more or add a figure about this.}
We formulate knowledge graph{\textendash}grounded dialog generation as follows:
\begin{equation}
\label{eq:problem}
    p_{\theta}(\OutSeq |\TokenSeq, \KG) = \prod_{j=1}^t p_{\theta}(\Out_j | \TokenSeq, \OutSeq_{<j}, \KG),
\end{equation}
where $\OutSeq=\left[y_1, y_2, \dots, y_t \right]$ is the output response, $t$ is the length of the response, and $\OutSeq_{<j}=\left[\Out_1, \dots \Out_{j-1} \right]$ denotes the generated sequence at the previous time steps. %, which is used for the autoregressive generation.
Since a KG can include a huge number of irrelevant entities and relations, KG-grounded dialog generation works generally retrieve subgraphs related to the dialog context for the efficiency and effectiveness.
% In the following sections, we describe how to retrieve and represent the subgraphs with a generative approach.
%Also, effectively representing the knowledge is crucial to reasoning over both the dialog and knowledge graph.